\chapter{RESULTS}
\label{ch:results}

\rev{The corrections in Chapters~\ref{ch:dmso} and~\ref{ch:control} change the
admissible gain ranges and, in the control case, the structure of the law
itself, so the results are regenerated here rather than carried over from
revision~7. All simulations and hardware replays in this chapter are produced by
the open-source Python implementation accompanying this thesis
(\texttt{twip-dmso}; script \texttt{scripts/thesis\_rev8.py}), which also records
every number quoted below.}

\section{TEST MATRIX AND METRICS}
\label{sec:testmatrix}

\subsection{Cases}

Each case is run for both the observer study and the control study.

\begin{table}[htbp]
\centering
\caption{Simulation test matrix}
\label{tab:testmatrix}
\begin{tabular}{clccc}
\toprule
Case & Description & Noise & Param.\ error & Nonlinearity \\
\midrule
O1 & Nominal & --- & --- & --- \\
O2 & Parameter uncertainty & --- & \checkmark & --- \\
O3 & Noise and uncertainty & \checkmark & \checkmark & --- \\
O4 & Unmodeled dynamics & \checkmark & \checkmark & --- \\
C1 & Control, nominal & \checkmark & --- & --- \\
C2 & Control, uncertainty & \checkmark & \checkmark & --- \\
C3 & Control, deadzone & \checkmark & \checkmark & deadzone \\
C4 & Control, deadzone + backlash & \checkmark & \checkmark & both \\
\bottomrule
\end{tabular}
\end{table}

\subsection{Simulation setup}
\label{sec:setup}

\rev{The truth model is the nonlinear model of Section~\ref{sec:simmodel},
integrated with fourth-order Runge--Kutta at one tenth of the
$\Delta t=\SI{10}{ms}$ loop period of the hardware. The estimators and
controllers use the nominal linearization discretized by \eqref{eq:zoh}.

\emph{Parameter error.} Cases O2--O4 and C2--C4 propagate the truth with
$M_p\times1.1$, $M_w\times1.5$, $k_m\times1.25$, $k_e\times0.95$, and $R\times1.2$;
the observer and controller retain the nominal values. This is the
perturbation of the revision~7 extra-control study with the torque-constant
error reduced from $\times2.5$ to $\times1.25$. With the corrected model of
Section~\ref{sec:simmodel}, the nominal LQR cannot stabilize the revision~7
perturbation even with perfect state knowledge, and a comparison of estimators
on a plant that no controller holds upright is not informative.

\emph{Unmodeled dynamics.} Case O4 adds the revision~7 unmodeled terms
$0.5\sin(x_1)x_2$ to $\dot x_2$ and $0.25x_2^2$ to $\dot x_4$.

\emph{Observer study.} For each case a single truth trajectory is generated by
the LQR of Section~\ref{sec:lqr} acting on the \emph{true} state, tracking a
position sinusoid of amplitude \SI{0.3}{m} and period \SI{8}{s} from an initial
tilt of \SI{5}{\degree}, for \SI{30}{s}. Every estimator then processes the same
measurement and input sequences, so all are compared on identical information.
The measurement $y_k$ is the full state as the robot constructs it: the encoder
position and velocity corrected for body tilt, the complementary-filtered tilt
of \eqref{eq:compfilter}, and the gyroscope rate. Without noise, $y_k=x_k$.
Each estimator is scored by its one-step prediction of $x_{k+1}$ after
processing $y_k$ and $u_k$, which is the native output of the DMSO and is
formed for the Kalman filters as $F\xhat_{k|k}+Gu_k$. The uncertainty truth is
$f_k=\Bm^T\bigl(x_{k+1}-Fx_k-Gu_k\bigr)$ computed from the true state.

\emph{Control study.} Each controller closes the loop itself for \SI{20}{s},
from \SI{5}{\degree} of tilt at rest, tracking a constant velocity of
\SI{0.1}{m/s}. The state is supplied to every controller by the same
Kalman filter, which models the accelerometer kinematics as described in
Section~\ref{sec:simmodel}. The LQR includes the nominal steady-state
feedforward voltage for the commanded velocity. Case C3 uses a
\SI{0.5}{V} deadzone and C4 adds an input backlash of half-width
$\Delta_b=\SI{0.25}{V}$.}

\subsection{Metrics}

\rev{Scalar summaries are reported per case rather than averaged across cases.
An average of percentage improvements taken over heterogeneous scenarios is not
a meaningful statistic, since the scenarios have different error scales and the
mean is dominated by whichever case happens to have the smallest baseline.}
The reported metrics are, for each state and each case:
\begin{enumerate}[label=(\roman*),leftmargin=*]
\item RMS estimation error over the run, after excluding the first
\SI{5}{s} as initial transient;
\item time to first entry into the ultimate bound predicted by
\eqref{eq:bounds}, together with that predicted bound;
\item RMS regulation or tracking error for the control cases;
\item RMS and peak control effort, to confirm that improvements are not bought
with unreported actuator authority.
\end{enumerate}

\subsection{Verification of the theoretical conditions}

\rev{Each run records $\sigmax(A)$, $\phi_{\max}$ as realized along the
trajectory, and the product $\Gamma\phi_{\max}^2$, so that satisfaction of
\eqref{eq:simple} can be confirmed rather than assumed.}

\rev{The DMSO is configured by the constructive procedure of
Appendix~\ref{app:validation}. The error-injection gain is $\Km=F-\bar\sigma I$,
so that $A=\bar\sigma I$ is normal and $\sigmax(A)=\bar\sigma=0.5$ exactly. The
basis is $\phi(y)=\bigl[1,\ \tanh(y_1/1),\ \tanh(y_2/1),\ \tanh(y_3/0.2),\
\tanh(y_4/1)\bigr]^T$, which is bounded by construction with
$\phi_{\max}^2=5$. The adaptation gain is one tenth of the bound
\eqref{eq:simple}, $\Gamma=0.01$. Table~\ref{tab:obscond} confirms that the
conditions hold along every trajectory; the realized $\phi_{\max}^2$ never
exceeds $2.2$, well below its design value.}

\begin{table}[htbp]
\centering
\caption{DMSO stability conditions as realized in each observer case}
\label{tab:obscond}
\small
\begin{tabular}{lccccccc}
\toprule
Case & $\sigmax(A)$ & $\phi_{\max}^2$ (design) & $\phi_{\max}^2$ (realized) & $\Gamma$
     & $\Gamma\phi_{\max}^2$ & bound on $\Gamma$ & \eqref{eq:simple} holds \\
\midrule
O1 & 0.50 & 5 & 1.86 & 0.0100 & 0.0186 & 0.100 & yes \\
O2 & 0.50 & 5 & 2.01 & 0.0100 & 0.0201 & 0.100 & yes \\
O3 & 0.50 & 5 & 2.19 & 0.0100 & 0.0219 & 0.100 & yes \\
O4 & 0.50 & 5 & 2.20 & 0.0100 & 0.0220 & 0.100 & yes \\
\bottomrule

\end{tabular}
\end{table}

\section{OBSERVER RESULTS}

\rev{Six estimators are compared: the three Kalman baselines of
Section~\ref{sec:dkf}; the revision~7 DMSO (basis on $\xhat_k$, $\Km=I$,
$\Gamma=0.1$) for reference; and the revision~8 DMSO, with and without
$\sigma$-modification ($\kappa=0.1$). A seventh row, marked $\dagger$, uses the
same adaptation law with the steady-state Kalman gain as error injection,
$\Km=FL$, so that $A=F(I-L)$. That $A$ is Schur but strongly non-normal
($\sigmax(A)=1.04$ without noise and $6.3$ with it), so Theorem~\ref{thm:main}
does not apply in the Euclidean norm and the row is an empirical extension;
see Remark~\ref{rem:nonnormal} and Section~\ref{sec:obsdiscussion}.}

\subsection{Nominal and parameter-uncertainty cases}

\rev{In case O1 there is no model error and every estimator predicts the state
to within numerical precision (Figure~\ref{fig:O1}); the DMSO weights remain
near zero because there is nothing to learn.

Case O2 is the setting the DMSO is designed for: a deterministic model error and
a noise-free measurement. Figure~\ref{fig:O2} shows the norm of the estimation
error against the ultimate bound of Theorem~\ref{thm:main}. The bound requires
$\epsilon_N$; a valid value on the realized trajectory is the largest residual
of a least-squares fit of the true uncertainty onto the basis,
$\epsilon_N=0.0113$, which gives $b_e=0.0235$. The error enters the bound after
0.12~s and remains inside it from 0.86~s onward. The excursions before that are
the initial transient, which Theorem~\ref{thm:main} does not constrain.

The DMSO identifies the uncertainty itself, not only the state
(Figure~\ref{fig:O2unc}). Against the augmented-state Kalman filter, the
baseline against which improvement is claimed, the RMS error of $\fhat$ is
lower by a factor of 11 in both the velocity and the tilt-rate channel
(Table~\ref{tab:obsunc}). The one-step prediction error of the uncertain states
is lower by a factor of 5 in both channels (Table~\ref{tab:obsrms}). Relative to
the revision~7 observer, the corrected basis and adaptation law reduce the
uncertainty estimation error by a factor of about 3. $\sigma$-modification
costs accuracy here, as the analysis predicts, because the leakage biases
the weights toward zero while the uncertainty is persistently nonzero.}

\begin{figure}[htbp]
\centering
\includegraphics[width=0.85\textwidth]{O1}
\caption{Case O1: states, DMSO one-step predictions, and prediction error.}
\label{fig:O1}
\end{figure}

\begin{figure}[htbp]
\centering
\includegraphics[width=0.85\textwidth]{O2_bound}
\caption{Case O2: norm of the estimation error with parameter uncertainty,
against the predicted ultimate bound $b_e$ of \eqref{eq:bounds}.}
\label{fig:O2}
\end{figure}

\begin{figure}[htbp]
\centering
\includegraphics[width=0.85\textwidth]{O2_uncertainty}
\caption{Case O2: DMSO uncertainty estimates against truth, with the revision~7
DMSO and the augmented-state Kalman filter for comparison.}
\label{fig:O2unc}
\end{figure}

\subsection{Cases with sensor noise}

\rev{With sensor noise the picture is mixed, and it is reported as such
(Figure~\ref{fig:O3} and Table~\ref{tab:obsrms}).

In the velocity channel the DMSO remains better than both unaugmented Kalman
filters, by a factor of 2.5 in case O3, and is within 40\% of the
augmented-state filter. In the tilt channel it is an order of magnitude worse
than the two unaugmented filters and comparable to the augmented one. The cause is the measurement rather than the adaptation.
The complementary-filtered tilt carries a slowly varying error of about
\SI{1.1}{\degree} RMS, because the accelerometer measures tilt relative to the
apparent vertical, which is displaced by the robot's own acceleration. With
$A=\bar\sigma I$ the DMSO corrects each channel only from its own measurement,
and so it follows that error. The Kalman filters instead infer tilt through
the model from the measured wheel acceleration, which the $A_2x_3$ coupling
makes available.

The Kalman-gain variant ($\dagger$) combines the two mechanisms. Its velocity
and tilt errors lie between those of the unaugmented and augmented filters, and
its tilt-rate error is 30\% above the nominal filter's in O3 and equal to it in
O4. It gives the best
uncertainty estimate in the noisy cases: 1.4 times lower
than the augmented filter in the velocity channel of O3, and 3 to 4 times lower
in the tilt-rate channel. Its gain lies outside the conditions of
Theorem~\ref{thm:main}, so these numbers are empirical.

In the tilt-rate channel of the noisy cases, every estimator's uncertainty error
is comparable to or larger than the RMS of the true uncertainty itself
(Table~\ref{tab:obsunc}). At this
noise level, $f_2$ is not identifiable from the available measurements by any
of the six estimators; only $f_1$ is.}

\begin{figure}[htbp]
\centering
\includegraphics[width=0.9\textwidth]{O3}
\caption{Case O3: estimation error (\SI{0.5}{s} moving RMS), DMSO against the
three Kalman baselines of Section~\ref{sec:dkf}.}
\label{fig:O3}
\end{figure}

\begin{table}[htbp]
\centering
\caption{Observer RMS one-step prediction error by case and estimator:
velocity (\si{mm/s}) / tilt (\si{mrad}) / tilt rate (\si{mrad/s}). Position
errors are set by the encoder quantization and are identical across
estimators (\SI{1.6}{mm} in O3 and O4).}
\label{tab:obsrms}
\scriptsize
\setlength{\tabcolsep}{3pt}
\begin{tabular}{lcccc}
\toprule
Estimator & O1 & O2 & O3 & O4 \\
\midrule
KF, nominal $Q_{\mathrm{K}}$ & $<$0.01 / $<$0.01 / $<$0.01 & 1.43 / 0.01 / 2.29 & 6.22 / 1.55 / 1.41 & 4.74 / 1.34 / 1.42 \\
KF, inflated $Q_{\mathrm{K}}$ & $<$0.01 / $<$0.01 / $<$0.01 & 1.43 / 0.01 / 2.29 & 6.63 / 3.81 / 1.10 & 6.42 / 3.68 / 1.12 \\
Augmented-state KF & $<$0.01 / $<$0.01 / $<$0.01 & 0.49 / 0.02 / 2.06 & 1.77 / 14.49 / 0.69 & 1.84 / 14.09 / 0.69 \\
DMSO, revision 7 & $<$0.01 / $<$0.01 / 0.01 & 0.14 / 0.01 / 0.54 & 6.17 / 20.09 / 1.25 & 6.08 / 19.48 / 1.21 \\
DMSO & $<$0.01 / $<$0.01 / 0.01 & 0.09 / 0.02 / 0.38 & 2.48 / 20.13 / 2.85 & 2.44 / 19.52 / 3.03 \\
DMSO with $\sigma$-mod & $<$0.01 / $<$0.01 / $<$0.01 & 0.61 / 0.02 / 0.88 & 2.71 / 20.13 / 4.91 & 2.60 / 19.52 / 5.10 \\
DMSO, Kalman-gain $\Km$\,$^\dagger$ & $<$0.01 / $<$0.01 / $<$0.01 & 0.96 / 0.06 / 1.47 & 1.97 / 6.51 / 1.86 & 2.05 / 3.95 / 1.39 \\
\bottomrule

\end{tabular}\\[2pt]
{\footnotesize $^\dagger$ outside the conditions of Theorem~\ref{thm:main}; empirical.}
\end{table}

\begin{table}[htbp]
\centering
\caption{RMS uncertainty estimation error $\fhat-f$ ($\times10^{-3}$),
velocity channel / tilt-rate channel. Only estimators that produce an
uncertainty estimate are listed; O1 has no uncertainty.}
\label{tab:obsunc}
\small
\begin{tabular}{lccc}
\toprule
Estimator & O2 & O3 & O4 \\
\midrule
Augmented-state KF & 0.69 / 11.27 & 0.59 / 7.28 & 0.59 / 7.25 \\
DMSO, revision 7 & 0.11 / 0.75 & 1.22 / 11.65 & 1.22 / 11.61 \\
DMSO & 0.44 / 2.81 & 0.94 / 11.44 & 0.94 / 11.36 \\
DMSO with $\sigma$-mod & 0.52 / 3.25 & 1.12 / 12.57 & 1.14 / 12.51 \\
DMSO, Kalman-gain $\Km$\,$^\dagger$ & 1.41 / 8.84 & 0.84 / 9.47 & 0.88 / 9.38 \\
\midrule
RMS of true $f$ & 1.99 / 12.38 & 1.99 / 12.38 & 1.98 / 12.24 \\
\bottomrule

\end{tabular}
\end{table}

\subsection{Discussion}
\label{sec:obsdiscussion}

\rev{Three conclusions follow. First, the corrected DMSO does what the analysis
claims. Its conditions are satisfied with margin, its error settles inside the
predicted bound, and under deterministic model error it identifies the
uncertainty an order of magnitude more accurately than the fairest Kalman
baseline. Second, the choice $A=\bar\sigma I$, which makes the proof simplest,
discards the cross-channel information that a model-based filter exploits. When
one measurement carries a structured error, that choice costs more than the
adaptation gains. Third, Theorem~\ref{thm:main} admits any $\Km$ with
$\sigmax(F-\Km)<1$, and the empirical success of the Kalman-gain injection
suggests extending the proof through the weighted norm of
Remark~\ref{rem:nonnormal}. That extension is not immediate. With
$V_k=e_k^TT^TTe_k+\Gamma^{-1}\norm{\Wt_k}_F^2$, the cancellations of
Proposition~\ref{prop:exact} survive only if the adaptation law is weighted
correspondingly, $\What_{k+1}=\What_k+\Gamma\phi_ke_{k+1}^TT^TT\Bm$, and if
$\Bm^TT^TT\Bm=I$. Those conditions should be imposed explicitly.}

\section{CONTROL RESULTS}

\subsection{Realizability of the command-filtered design}
\label{sec:cfinstab}

\rev{The command-filtered design of Section~\ref{sec:cffb} could not be
stabilized, and the reason is structural. Stage~3 of
Appendix~\ref{app:validation} was followed: the backstepping structure was run
first with both networks disabled, on the nominal linear model, with perfect
state knowledge and no saturation. The closed loop diverged within seven samples.

Linearizing the continuous-time closed loop (plant and three command filters,
networks disabled) gives a real eigenvalue near $+7$~rad/s for every gain set
tried: 720 combinations of $c_1,\dots,c_4$ and $\varpi_1,\varpi_2,\varpi_3$, including
all combinations satisfying \eqref{eq:gaincond}. The smallest value found for
the largest real part was $+6.35$~rad/s. As the filter bandwidths are raised,
that eigenvalue converges to $+7.23$~rad/s (Figure~\ref{fig:chi}), which is the
right-half-plane transmission zero of the position-to-voltage transfer function
$x_1(s)/u(s)$. On the revision~7 parameters the zero is at $+6.11$~rad/s, so the
difficulty does not come from the model corrections.

The mechanism is the following. With $x_1$ as the tracked output, the recursion
drives all four error coordinates to zero, which requires the closed loop to
invert the plant from $u$ to $x_1$. For a non-minimum-phase output, that
inversion places a closed-loop pole at the zero. In the analysis the
instability is hidden in Assumption~\ref{as:filt}. The virtual control
$\alpha_2$ contains $B_1u$, the control depends on the filter state
$\dot x_4^c$, and $\alpha_3$ contains $\dot x_3^c$. So $\alpha_2$ and its
derivatives depend on the filter errors themselves, and the premise that
$\alpha_i$ and $\dot\alpha_i$ are bounded, which yields
$\chi_i=O(1/\varpi_i)$, does not follow from the construction. Raising the
bandwidth moves the unstable eigenvalue toward the zero rather than into the
left half-plane.

Two remedies preserve the intent of Section~\ref{sec:cffb}. The first is to
track a minimum-phase output. For a controllable single-input linearization, the
flat output $y=Cx$ with $CB=CAB=CA^2B=0$ has relative degree four and no zero
dynamics, and a four-step recursion on it is well posed. The second is a
cascade: a position loop slower than the zero that commands tilt, and a
backstepping tilt loop with the network compensation. The control table below
therefore reports the command-filtered design as unrealizable in its present
form.}

\begin{figure}[htbp]
\centering
\includegraphics[width=0.8\textwidth]{cfb_eigenvalues}
\caption{Largest real part of the closed-loop eigenvalues of the command-filtered
design (networks disabled, nominal linear model) against the command-filter
bandwidth, for three sets of damping gains. The dashed line is the
right-half-plane zero of $x_1(s)/u(s)$. Assumption~\ref{as:filt} requires the
filter errors $\chi_i$ to shrink as $\varpi_i$ grows; instead the unstable
eigenvalue converges to the zero.}
\label{fig:chi}
\end{figure}

\subsection{LQR and the two-step extra control}

\rev{Table~\ref{tab:ctrleff} and Figures~\ref{fig:C2}--\ref{fig:C4} compare
the LQR alone with the revision~7 two-step extra control. The two-step design
is the \texttt{ExtraControl\_v5.m} law with $k_{e0}=0$, the setting that
reproduces the revision~7 figures.

The two-step extra control reduces the RMS position tracking error by 17\% in
the nominal case C1 and by 30\% under deadzone and backlash (C4). It increases
it by 4\% with parameter uncertainty alone (C2) and by 22\% with the deadzone
(C3). It uses between 1\% and 29\% more RMS control. Both controllers balance in
every case. In both, the position error is dominated by the slow position
mode of the LQR (the thesis weights place a closed-loop pole at
$|\lambda|=0.9995$ at \SI{10}{ms}). Under the deadzone it is also dominated by a
steady velocity deficit that neither design removes, visible as the ramp in
Figure~\ref{fig:C4}.

A slow position ramp, the behavior Remark~\ref{rem:twostep} predicts for the
two-step design, appears in cases C2--C4. It appears equally for the LQR alone,
with slopes of 0.02 to \SI{0.06}{m/s} for both controllers over the last
\SI{10}{s}. So it is not a property of the two-step design. It is the
velocity error left by an incorrect steady-state feedforward under parameter
error, integrated by a position loop too weak to remove it. Test T3 of
Appendix~\ref{app:validation} therefore neither confirms nor refutes
Remark~\ref{rem:twostep}. The LQR term supplies whatever position damping
there is in both designs.}

\begin{figure}[htbp]
\centering
\includegraphics[width=0.85\textwidth]{C2}
\caption{Case C2: position and velocity tracking errors, tilt, and tilt rate
with and without the two-step extra control.}
\label{fig:C2}
\end{figure}

\begin{figure}[htbp]
\centering
\includegraphics[width=0.85\textwidth]{C4}
\caption{Case C4: response under deadzone and backlash.}
\label{fig:C4}
\end{figure}

\begin{table}[htbp]
\centering
\caption{Control effectiveness by case: RMS position tracking error (\si{m})
after \SI{5}{s} / RMS voltage (\si{V}) / peak voltage (\si{V}).}
\label{tab:ctrleff}
\footnotesize
\setlength{\tabcolsep}{4pt}
\begin{tabular}{lcccc}
\toprule
Controller & C1 & C2 & C3 & C4 \\
\midrule
LQR alone & 0.055 / 0.72 / 2.6 & 0.068 / 0.73 / 2.6 & 0.091 / 0.74 / 2.6 & 0.680 / 0.39 / 2.3 \\
LQR + two-step (rev.~7) & 0.051 / 0.72 / 2.6 & 0.068 / 0.73 / 2.6 & 0.100 / 0.75 / 2.6 & 0.647 / 0.41 / 2.3 \\
Command-filtered, flat output, networks off & 0.004 / 0.93 / 9.7 & 0.016 / 0.90 / 9.6 & 0.016 / 0.91 / 9.7 & 0.033 / 0.93 / 9.6 \\
Command-filtered, flat output, networks on & 0.004 / 0.93 / 9.7 & 0.016 / 0.90 / 9.6 & 0.016 / 0.91 / 9.7 & 0.031 / 0.93 / 9.6 \\
Command-filtered, as in Section~\ref{sec:cffb} & \multicolumn{4}{c}{unstable for all gains tested (Section~\ref{sec:cfinstab})} \\
\bottomrule

\end{tabular}
\end{table}

\section{IMPLEMENTATION RESULTS}
\label{sec:implresults}

\rev{The hardware results are replays of the logged LQR run
\texttt{implementationtest5.txt}, recorded with firmware v9 at a
\SI{10}{ms} loop. Figure~\ref{fig:hw} compares the measured states, the
estimates computed on board by the v9 DMSO (which the replay reproduces to
print precision), and the revision~8 DMSO run offline on the same data. After
the first second the three agree to within the measurement quantization. There
is no truth on hardware, so the replays establish consistency, not accuracy.}

\begin{figure}[htbp]
\centering
\includegraphics[width=0.85\textwidth]{hw_states}
\caption{Hardware test: measured states, on-board (v9) DMSO estimates, and the
revision~8 DMSO replayed offline.}
\label{fig:hw}
\end{figure}

\begin{figure}[htbp]
\centering
\includegraphics[width=0.85\textwidth]{hw_complementary}
\caption{Hardware test: complementary filter against the model-based estimators
using raw accelerometer tilt.}
\label{fig:complementary}
\end{figure}

\begin{figure}[htbp]
\centering
\includegraphics[width=0.85\textwidth]{hw_uncertainty}
\caption{Hardware test: DMSO uncertainty estimates.}
\label{fig:hwunc}
\end{figure}

\rev{The revision~7 result that the complementary filter outperforms the
model-based estimators fed raw accelerometer tilt is re-examined in
Figure~\ref{fig:complementary}. Without truth, each tilt estimate is scored
by its consistency with the gyroscope, the RMS of
$\dd\hat\theta/\dd t-g_y$. This measures high-frequency error and is blind to
slow bias. On that score the raw accelerometer tilt is at
\SI{476}{\degree/s}, the complementary filter at \SI{9.9}{\degree/s}, and the
DMSO fed raw accelerometer tilt at \SI{295}{\degree/s}. The Kalman filter fed
the same raw tilt, with a measurement variance that reflects its actual error,
reaches \SI{4.2}{\degree/s}.

The finding therefore holds for the DMSO but not for a correctly weighted
Kalman filter. The explanation needs no gearbox transients: an accelerometer
measures tilt relative to the apparent vertical, which is displaced by
$\ddot x/g$ whenever the robot accelerates. An estimator that corrects each
channel from its own measurement follows that displacement, and one that
weights the measurement by its true error does not. The conclusion of
revision~7 stands in modified form: it is a statement about measurement
construction and error-injection design, not about the adaptation.

Figure~\ref{fig:hwunc} shows the uncertainty estimated on hardware. The
tilt-rate channel is several times larger than in simulation (RMS
$1.4\times10^{-2}$ against $2.3\times10^{-3}$), consistent with the model
error on the physical robot being larger than the simulated perturbation.
Appendix~\ref{app:validation} describes how to settle this by measurement.}
